\section{Paired binary states and two factorial lifts}
\label{sec:signatures}

We now encode the small-cell counts by subsets of a polynomial-size
binary ground set.  The intended subsets select one
element from each labeled pair.  Several such transversals represent
the same balanced count profile, so a factorial normalization is needed
to recover $C(j)$ as their total mass.  Two normalizations agree on transversals but have
different quadratic signatures.  The first will support transport;
the second will compare the masses of the defect classes through
which the transport passes.

\subsection{The paired encoding and its weights}

Assume in the rest of this section that \(p\ge1\); the case \(p=0\)
will be handled directly by the weight evaluator.
For each small cell \(s\), introduce \(b_s\) labeled pairs, each
with a row element and a column element.
Number the pairs \(1,\ldots,p\), and let \(E\) be their \(2p\)
elements.
For a \(p\)-subset \(S\subseteq E\), write \(x_s(S),y_s(S)\) for
its row and column counts in cell \(s\).
For \(0\le j\le H\), define positive weights on \emph{all}
\(p\)-subsets by
\begin{align}
 f_j(S)
 &=h_j(x,y)\prod_{s\ \mathrm{small}}\frac{x_s!\,y_s!}{b_s!},
 \notag\\
 \widetilde f_j(S)
 &=h_j(x,y)\prod_{s\ \mathrm{small}}
       \frac{(b_s-x_s)!\,(b_s-y_s)!}{b_s!}.
 \label{eq:lifts}
\end{align}
All factorial arguments lie between zero and \(b_s\).
We suppress \(j\) when it is fixed.

Let \(\mathcal A\) be the transversals, consisting of exactly one
element from each labeled pair.
For distinct pair labels \(i,l\), let \(\mathcal A_{il}\) consist
of sets with pair \(i\) empty, pair \(l\) full, and every other
pair singly occupied.
On a transversal \(x_s+y_s=b_s\), so both lifts have the same
factorial multiplier \(1/\binom{b_s}{x_s}\) in cell \(s\).
There are exactly \(\prod_s\binom{b_s}{x_s}\) transversals with a
given balanced count profile.
Consequently
\begin{equation}
 \sum_{S\in\mathcal A}f(S)
 =\sum_{S\in\mathcal A}\widetilde f(S)
 =C(j).
 \label{eq:transversal-mass}
\end{equation}

For example, in a cell of bound three, the balanced counts $(2,1)$
are represented by the three choices of which labeled pair uses its
column element.  Each transversal has factorial factor
$2!1!/3!=1/3$, so their total contribution is exactly the original
count weight.  This cancels the multiplicity among transversals; it does not count
all subsets with the same coordinate counts.  For instance, these
counts also occur in subsets with one empty pair and one full pair.

Put \(q=W^2\).
On the enlarged state space used later,
\begin{equation}
 q^{-1}f(S)\le\widetilde f(S)\le qf(S),
 \qquad
 S\in\mathcal A\cup\bigcup_{i\ne l}\mathcal A_{il}.
 \label{eq:lift-comparison}
\end{equation}
To see this, first suppose the empty and full pairs lie in the
same cell.  Every cell remains balanced, and the ratio is one.
Otherwise one cell has count sum \(b_s-1\), contributing
\((x_s+1)(y_s+1)\) to \(\widetilde f/f\), and another has count
sum \(b_s+1\), contributing \(1/(x_sy_s)\).
Both positive integer products in this quotient lie between one
and \(W^2\), since the small bounds are less than \(W\).
Their quotient is therefore in \([W^{-2},W^2]\).
We do not assert \eqref{eq:lift-comparison} outside this enlarged
state space; the signature properties below, in contrast, hold
on the full \(p\)-subset domain.

\subsection{Binary signatures and their closure properties}

For a positive weight \(F\) on all \(k\)-subsets of a finite set,
the \emph{addition property} means the following.
After fixing any \(k-2\) included elements, the symmetric matrix
whose off-diagonal \(i,h\) entry is the weight obtained by including
the two further distinct elements \(i,h\), and whose diagonal is
zero, has at most one positive eigenvalue.
For \(k<2\) the property is vacuous.
The \emph{omission property} is the addition property of the
complement-set weight.

For a positive weight on every $k$-subset, the addition property is
the Lorentzian criterion for the multiaffine generating polynomial
$\sum_S F(S)\prod_{i\in S}z_i$; the omission property applies to the
complement-set polynomial~\cite[Theorem~2.25]{BH2020}.
Summing a pair $\{a,b\}$ at occupancy one corresponds to applying
$\partial_a+\partial_b$ and then setting $z_a=z_b=0$.
Both operations preserve the Lorentzian property by
\cite[Theorem~2.10 and Corollary~2.11]{BH2020}.
The matrix proof below also keeps track of the full index set before
that specialization, as required for the three-pair defect inequality.

\begin{lemma}[Binary contractions]
\label{lem:binary-contraction}
The addition property, with positivity on the entire valid-degree
domain, is preserved by fixing elements included or excluded and by
summing a two-element group at occupancy one.
The same assertions hold for the omission property.
In addition, the cubic contraction used to sum a pair has at most
one positive eigenvalue on its full remaining index set, before the
contracted pair's indices are removed.
\end{lemma}

\begin{proof}
Fixing inclusions enlarges the set of fixed elements in the
quadratic test, and fixing exclusions restricts that test to a
principal submatrix.
Positivity is retained whenever the residual degree is valid.

We prove the contraction assertion through a symmetric cubic
tensor.
Suppose \(T_{iuh}\) is positive on triples of distinct indices and
zero on repetitions, and every slice \(T_i=(T_{iuh})_{u,h}\)
has at most one positive eigenvalue.
There are at least three indices.
For a vector \(a\) with all entries positive, set
\[
 M=\sum_i a_iT_i.
\]
Tensor symmetry gives
\[
 a^{\mathsf T}T_i a=(Ma)_i>0,
 \qquad
 h^{\mathsf T}T_i a=(Mh)_i.
\]
The strict inequality holds even with exactly three indices:
the two indices other than \(i\) give a strictly positive term.
Applying \eqref{eq:signature-projection} to every slice and then
summing with coefficients \(a_i\) yields
\begin{equation}
 h^{\mathsf T}Mh
 \le \sum_i\frac{a_i(Mh)_i^2}{(Ma)_i}
 =h^{\mathsf T}MD_*^{-1}Mh,
 \qquad
 D_*=\operatorname{diag}\left(\frac{(Ma)_i}{a_i}\right).
 \label{eq:cubic-projection}
\end{equation}
Thus the symmetric matrix \(M'=D_*^{-1/2}MD_*^{-1/2}\)
satisfies \(M'\preceq(M')^2\).
It has the positive eigenvector \(u=D_*^{1/2}a\) with eigenvalue
one.
Every off-diagonal entry of \(M\), and hence of \(M'\), is
strictly positive, because a third distinct index contributes
a positive term.
The matrix with entries \(M'_{ih}u_h/u_i\) is therefore
irreducible and row-stochastic.
The maximum principle bounds the absolute values of its real
eigenvalues by one and shows that its eigenvalue-one eigenspace
consists of the constant vectors.
Since \(M'\) is symmetric, eigenvalue one is simple.
On the other hand, \eqref{eq:cubic-projection} requires every
eigenvalue \(\lambda\) of \(M'\) to satisfy
\(\lambda\le\lambda^2\), so a positive eigenvalue must be at least
one.
There can consequently be only one positive eigenvalue.
Positive diagonal congruence gives the same conclusion for \(M\).
For nonnegative \(a\), replace it by \(a+\varepsilon\mathbf 1\)
and pass to the matrix limit as \(\varepsilon\downarrow0\).
This establishes the full-matrix contraction assertion.

Now sum a pair \(\{\alpha,\beta\}\) at occupancy one in a positive
weight \(F\) on \(k\)-subsets:
\[
 F^\downarrow(U)=F(U\cup\{\alpha\})+F(U\cup\{\beta\}),
 \qquad
 |U|=k-1,\quad
 U\cap\{\alpha,\beta\}=\varnothing.
\]
If \(k-1<2\), the desired addition property is vacuous.
Otherwise fix \(k-3\) included elements outside the pair.
The remaining three-element weights form a tensor \(T\) of the
kind just considered.
Each slice is an addition matrix of \(F\), with a zero row and
column at its fixed index.
Contract with the indicator of \(\{\alpha,\beta\}\).
Restricting the resulting full matrix to the other indices gives
exactly the addition matrix of \(F^\downarrow\).
The result is therefore valid.
Every valid residual weight is a sum of positive terms, so
positivity also persists.

Finally, complementation interchanges fixed inclusions with fixed
exclusions and preserves occupancy one in a two-element pair.
Applying the addition argument to complement weights proves all
the omission assertions.
\end{proof}

\begin{lemma}[Signatures of the two lifts]
\label{lem:binary-signatures}
The weight \(f_j\) has the omission property, and
\(\widetilde f_j\) has the addition property.
These properties and positivity remain valid after fixing any
included or excluded elements and summing any disjoint labeled
pairs at occupancy one, whenever the residual degree is valid.
All assertions concern the full valid-degree subset domain.
\end{lemma}

\begin{proof}
The case \(p=1\) is immediate for the original two properties, so
consider a quadratic test when it is nonvacuous.
For the omission property of \(f\), fix \(p-2\) omitted elements.
This leaves a display of \(p+2\) elements from which two must be
removed.
Group its elements by small-cell row or column coordinate, and
write \(n_u\) for the number in group \(u\).

Within a group, all distinct-copy matrix entries are equal, and
the diagonal entries are zero.
Vectors whose entries sum to zero in each group have nonpositive
square: in each group the form is a nonnegative multiple of
\(\mathbf 1\mathbf 1^{\mathsf T}-I\).
These vectors have no coupling to vectors that are constant within
each group.
It remains to examine the latter subspace.

Let its value on group \(u\) be \(t_u\).
For distinct groups \(u,v\), the \(n_un_v\) ordered-copy
multiplicity cancels the two factorial reductions when passing
from displayed counts \(n\) to \(n-e_u-e_v\).
For a single group \(u\), the multiplicity \(n_u(n_u-1)\)
similarly cancels the factorial reduction for two removals.
If \(n_u<2\) this term is zero.
The resulting quadratic form is exactly
\[
 \frac{\prod_u n_u!}{\prod_s b_s!}
 \sum_{u,v}
 h_j(n-e_u-e_v)t_ut_v,
\]
with zero entries for invalid removals.
By Corollary~\ref{cor:hj-signatures}, it has at most one positive
eigenvalue.
Together with the nonpositive group-sum-zero subspace, this proves
the omission property of \(f\).

For the addition property of \(\widetilde f\), fix \(p-2\)
included elements, with coordinate counts \(r_u\).
There are \(n_u=b_u-r_u\) available elements in coordinate group
\(u\), where \(b_u=b_s\) for either coordinate of cell \(s\).
After adding two elements, the factorial numerator of
\(\widetilde f\) is the product of the remaining-capacity
factorials.
The same group decomposition and the same multiplicity cancellation
therefore give, up to a common positive scalar, the form
\[
 \sum_{u,v}h_j(r+e_u+e_v)t_ut_v.
\]
Invalid additions give zero entries.
The addition assertion of Corollary~\ref{cor:hj-signatures}
proves that this form has at most one positive eigenvalue.
The group-sum-zero subspace again contributes only nonpositive
squares.
This proves the addition property of \(\widetilde f\).

All \(p\)-subset weights in \eqref{eq:lifts} are strictly positive.
The closure assertions now follow from
Lemma~\ref{lem:binary-contraction}, applied respectively in omission
and addition variables.
\end{proof}

\begin{remark}[The two signature roles are distinct]
\label{rem:distinct-lifts}
Even with a constant count weight, the factorial construction need
not give \(f\) the addition property.
For example, take two small cells with bounds \(3\) and \(2\),
with base weight \(h\equiv1\), so the binary ground set has ten
elements and \(p=5\).
Fix as included the three row elements of the first cell.
The seven remaining elements are its three column elements and
the four elements of the second cell.
In the addition matrix of \(f\), restrict to vectors with common
value \(a\) on the first group of three and common value \(b\) on
the second group of four.
Direct substitution in \eqref{eq:lifts} gives the quadratic form
\[
 6a^2+12ab+8b^2.
\]
Its coefficient matrix
\(\left(\begin{smallmatrix}6&6\\6&8\end{smallmatrix}\right)\)
is positive definite, with determinant \(12\).
Thus this addition matrix has at least two positive eigenvalues.
The omission property of \(f\) and the addition property of
\(\widetilde f\) are the respective assertions needed below.
\end{remark}

Figure~\ref{fig:paired-states} illustrates the paired state space,
the agreement of the lifts on transversals, and their distinct
signature roles.
\begin{figure}[tbp]
  \centering
  \begin{tikzpicture}[
    x=1cm,y=1cm,
    font=\small,
    element/.style={circle,draw=black!80,line width=.6pt,
      minimum size=4.7mm,inner sep=0pt},
    selected/.style={element,fill=black!80},
    pair/.style={draw=black!22,rounded corners=2pt,line width=.5pt},
    lift/.style={draw=black!35,rounded corners=2pt,
      fill=black!2,align=center,inner sep=7pt,
      text width=4.65cm,minimum height=1.55cm}
  ]
    \node[font=\small\bfseries] at (2.2,2.52) {Transversal $S\in\mathcal A$};
    \node[font=\small\bfseries] at (8.1,2.52) {One-defect state $S'\in\mathcal A_{il}$};

    \foreach \x/\lab in {1.15/i,1.90/k,2.65/l,3.40/m} {
      \draw[pair] (\x-.27,.43) rectangle (\x+.27,1.72);
      \node at (\x,1.99) {$\lab$};
      \node[element] at (\x,1.41) {};
      \node[element] at (\x,.74) {};
    }
    \node[anchor=east] at (.73,1.41) {row};
    \node[anchor=east] at (.73,.74) {column};
    \foreach \x/\y in {1.15/1.41,1.90/1.41,2.65/.74,3.40/.74}
      \node[selected] at (\x,\y) {};
    \node at (2.2,.11) {one selected per pair};
    \node at (2.2,-.40) {$f(S)=\widetilde f(S)$};

    \foreach \x/\lab in {7.05/i,7.80/k,8.55/l,9.30/m} {
      \draw[pair] (\x-.27,.43) rectangle (\x+.27,1.72);
      \node at (\x,1.99) {$\lab$};
      \node[element] at (\x,1.41) {};
      \node[element] at (\x,.74) {};
    }
    \foreach \x/\y in {7.80/1.41,8.55/1.41,8.55/.74,9.30/.74}
      \node[selected] at (\x,\y) {};
    \draw[draw=blue!55!black,dashed,rounded corners=2pt,line width=.8pt]
      (6.74,.39) rectangle (7.36,1.76);
    \draw[draw=red!55!black,double,double distance=.7pt,
      rounded corners=2pt,line width=.45pt]
      (8.24,.39) rectangle (8.86,1.76);
    \node[font=\footnotesize,text=blue!55!black] at (7.05,.12) {empty};
    \node[font=\footnotesize,text=red!55!black] at (8.55,.12) {full};
    \node at (8.1,-.40) {$q^{-1}f(S')\leq\widetilde f(S')\leq qf(S')$};
    \node[font=\footnotesize] at (8.1,-.80) {$q=W^2$};

    \draw[-{Stealth[length=2mm]},line width=.7pt]
      (4.00,1.10) -- (5.65,1.10);
    \node[font=\footnotesize,align=center] at (4.825,1.65)
      {one exchange};
    \node[font=\footnotesize,align=center] at (4.825,.53)
      {remove row $i$\\add row $l$};

    \node[lift] at (2.35,-2.00)
      {$f$: \textbf{omission signature}\\[3pt]
       transport and simulation};
    \node[lift] at (8.00,-2.00)
      {$\widetilde f$: \textbf{addition signature}\\[3pt]
       inequalities for defect masses};
    \node[font=\footnotesize] at (5.175,-3.02)
      {Both signature properties are proved on all $p$-subsets.};
    \node[selected,minimum size=3.2mm] at (3.81,-3.49) {};
    \node[anchor=west,font=\footnotesize] at (4.02,-3.49) {selected};
    \node[element,minimum size=3.2mm] at (5.73,-3.49) {};
    \node[anchor=west,font=\footnotesize] at (5.94,-3.49) {unselected};
  \end{tikzpicture}
  \caption{The paired state space and the two factorial lifts, illustrated with
  four pairs. The lifts agree on transversals. Their comparison within a factor
  $q$ is used only on transversals and one-defect types. The distinct signature
  properties support the two different parts of the argument shown below the
  states.}
  \label{fig:paired-states}
\end{figure}


\subsection{Defect masses and pair inequalities}

Fix \(j\), put \(C=C(j)\), and suppress the parameter on both lifts.
A partial transversal assignment \(\sigma\) specifies one chosen
element in each of a collection of labeled pairs.
A set respects \(\sigma\) when its intersection with every assigned
pair is precisely that chosen element.
For unassigned distinct pair labels \(i,l\), define
\begin{align*}
 z_\sigma
 &=\sum_{\substack{S\in\mathcal A\\S\text{ respects }\sigma}}f(S)
   =\sum_{\substack{S\in\mathcal A\\S\text{ respects }\sigma}}
       \widetilde f(S),\\
 c_{il}^{\sigma}
 &=\sum_{\substack{S\in\mathcal A_{il}\\
                   S\text{ respects }\sigma}}f(S),
 \qquad
 \widetilde c_{il}^{\sigma}
 =\sum_{\substack{S\in\mathcal A_{il}\\
                  S\text{ respects }\sigma}}\widetilde f(S).
\end{align*}
These masses are positive.
For the empty assignment we write \(c_{il},\widetilde c_{il}\)
without a superscript and have \(z_\varnothing=C\).

\begin{lemma}[Pair inequalities for defect masses]
\label{lem:defect-inequalities}
For every partial assignment \(\sigma\) and distinct unassigned
labels \(i,l\),
\begin{equation}
 4\widetilde c_{il}^{\sigma}\widetilde c_{li}^{\sigma}
 \le z_\sigma^2.
 \label{eq:defect-two-pair}
\end{equation}
For three distinct unassigned labels \(i,l,k\),
\begin{equation}
 \widetilde c_{il}^{\sigma}\widetilde c_{lk}^{\sigma}
 \le z_\sigma\widetilde c_{ik}^{\sigma}.
 \label{eq:defect-three-pair}
\end{equation}
\end{lemma}

\begin{proof}
Fix the inclusions and exclusions specified by \(\sigma\).
For \eqref{eq:defect-two-pair}, sum every other unassigned pair at
occupancy one, leaving only \(i,l\).
By Lemma~\ref{lem:binary-signatures}, the remaining positive
two-element weight has an addition matrix with at most one
positive eigenvalue.
Restrict its quadratic form to vectors constant on each of the two
pairs.
The resulting matrix is
\[
 \begin{pmatrix}
  2\widetilde c_{li}^{\sigma}&z_\sigma\\
  z_\sigma&2\widetilde c_{il}^{\sigma}
 \end{pmatrix}.
\]
The diagonal factor two counts the two orders of the elements in
a full pair; the off-diagonal sums exactly the transversals.
The matrix has positive trace and at most one positive eigenvalue,
so its determinant is nonpositive.
This proves \eqref{eq:defect-two-pair}.

For \eqref{eq:defect-three-pair}, sum every other unassigned pair
at occupancy one, leaving the six elements in \(i,l,k\) and a
positive weight on their three-element subsets.
Form its symmetric cubic tensor, zero on repeated indices.
Contract it along the indicator of the two elements of pair \(k\).
By the full-matrix assertion of
Lemma~\ref{lem:binary-contraction}, the resulting matrix has at
most one positive eigenvalue even while the two contracted indices
are retained.
Restrict its form to vectors constant on each pair, and call the
resulting three-by-three matrix \(N\), indexed by \(i,l,k\).
Its relevant entries are
\begin{equation}
 N_{ll}=2\widetilde c_{il}^{\sigma},
 \qquad
 N_{ik}=2\widetilde c_{lk}^{\sigma},
 \qquad
 N_{il}=z_\sigma,
 \qquad
 N_{lk}=2\widetilde c_{ik}^{\sigma}.
 \label{eq:contracted-defect-entries}
\end{equation}
For example, in \(N_{ik}\) the full pair \(k\) can contribute either
of its elements to the contraction and the other to the matrix
index, explaining the factor two.
All entries of \(N\) are nonnegative, and \(N_{ll}>0\).

By \eqref{eq:signature-projection}, the Schur complement against
the \(l,l\) entry is negative semidefinite.
Its two remaining coordinates \(i,k\) therefore satisfy
\begin{align*}
 N_{ik}
 &\le \frac{N_{il}N_{lk}}{N_{ll}}
 +\sqrt{
   \left(\frac{N_{il}^2}{N_{ll}}-N_{ii}\right)
   \left(\frac{N_{lk}^2}{N_{ll}}-N_{kk}\right)}\\
 &\le \frac{2N_{il}N_{lk}}{N_{ll}}.
\end{align*}
The square-root factors are nonnegative by the same Schur
complement condition.
The second inequality uses \(N_{ii},N_{kk}\ge0\) and the
nonnegativity of the off-diagonal entries.
Substituting \eqref{eq:contracted-defect-entries} gives
\eqref{eq:defect-three-pair}.
\end{proof}

The omission property of $f$ is the input to the one-hole transport
argument in Section~\ref{sec:transport}; it applies even to auxiliary
completions outside the enlarged state space.  The two addition-based
inequalities just proved control the conditional and unconditional
one-defect masses used to bound transversal variance and defect means.
The comparison of the lifts is used only on the actual transversal
and one-defect classes where \eqref{eq:lift-comparison} holds.

The matrices and balanced sums used to prove the signature
properties are analytic devices.
The algorithm will neither form these matrices nor enumerate their
underlying coordinate boxes.
